% Part: normal-modal-logic
% Chapter: syntax-and-semantics
% Section: relational-models

\documentclass[../../../include/open-logic-section]{subfiles}

\begin{document}

\olfileid{nml}{syn}{rel}

\olsection{संबंधी प्रतिमाने}

सामान्य मोडल तर्कशास्त्रांच्या चिन्हार्थमीमांसेची मूलभूत संकल्पना
म्हणजे \emph{संबंधी प्रतिमान}. त्यात द्विस्थानी ``प्राप्यता संबंधाने''
संबंधित जगांचा संच आणि कोणते !!{propositional variable}s कोणत्या
जगांत ``सत्य'' मानायची हे ठरवणारे मूल्यांकन असते.

\begin{defn}
  मूलभूत मोडल भाषेचे \emph{प्रतिमान} म्हणजे त्रिकूट $\mModel{M}
  = \tuple{W, R, V}$; येथे
  \begin{enumerate}
  \item $W$ हा ``जगांचा'' रिकामा नसलेला संच आहे,
  \item $R$ हा~$W$ वरील द्विस्थानी प्राप्यता संबंध आहे, आणि
  \item $V$ हे प्रत्येक !!{propositional
    variable}~$p$ ला संभाव्य जगांचा संच~$V(p)$ नेमून देणारे फलन आहे.
  \end{enumerate}
  $Rww'$ असेल तेव्हा $w'$ हे~$w$ पासून \emph{प्राप्य}
  आहे असे म्हणतो. $w \in V(p)$ असेल तेव्हा $p$ हे~$w$ येथे
  \emph{सत्य} आहे असे म्हणतो.
\end{defn}

संबंधी चिन्हार्थमीमांसेचा मोठा फायदा असा की प्रतिमाने
\olref{fig:simple} मधील साध्या आकृतीसारख्या आकृत्यांनी दाखवता
येतात. जगांसाठी गाठी वापरतात; $w$ पासून~$w'$ कडे बाण
असेल तेव्हा आणि केवळ तेव्हाच जग~$w'$ हे~$w$ पासून प्राप्य
असते. शिवाय $w \in V(p)$ असेल तर गाठीला (जगाला)~$\mTrue{p}$,
अन्यथा~\mFalse{p} अशी खूण करतो. \olref{fig:simple} मध्ये
$W = \{w_1, w_2, w_3\}$, $R = \{\tuple{w_1, w_2},\tuple{w_1,
  w_3}\}$, $V(p) = \{w_1, w_2\}$ आणि $V(q) = \{w_2\}$
असलेले प्रतिमान दाखवले आहे.

\begin{figure}
  \begin{center}
    \begin{tikzpicture}[modal]
      \node[world] (w1) [label={[align=right]right:\mTrue{p}\\ \mFalse{q}}]{$w_1$};
      \node[world] (w2) [label={[align=right]right:\mTrue{p}\\ \mTrue{q}},
        above right=of w1]{$w_2$};
      \node[world] (w3) [label={[align=right]right:\mFalse{p}\\ \mFalse{q}},
        below right=of w1] {$w_3$};
      \draw[->] (w1) to (w2);
      \draw[->] (w1) to (w3);
    \end{tikzpicture}
  \end{center}
  \caption{साधे प्रतिमान.}
  \ollabel{fig:simple}
\end{figure}

\end{document}
